\documentclass[tikz,border=8pt]{standalone}
\usepackage{tikz}
\usetikzlibrary{calc,arrows.meta}
\usepackage{amsmath}
\usepackage{pgfplots}
\pgfplotsset{compat=1.18}
\usepackage{ctex}

\begin{document}
\begin{tikzpicture}

% ---- 左图：先验分布 Beta(2,2) ----
\begin{axis}[
  name=ax1,
  width=5.5cm, height=5.5cm,
  xmin=-0.05, xmax=1.05,
  ymin=-0.1, ymax=2.5,
  xlabel={$\theta$},
  ylabel={密度},
  title={\textbf{先验} $\pi(\theta)$},
  axis lines=left,
  xtick={0,0.25,0.5,0.75,1},
  xticklabel style={font=\scriptsize},
  yticklabel style={font=\scriptsize},
  ytick={0,0.5,1.0,1.5,2.0},
  grid=major, grid style={gray!20},
  clip=false
]
  % Beta(2,2): f(x) = 6*x*(1-x)，峰值在 x=0.5，f(0.5)=1.5
  \addplot[blue!70!black, ultra thick, domain=0.01:0.99, samples=80]
    {6*x*(1-x)};
  \addplot[fill=blue!20, opacity=0.5, draw=none, domain=0.01:0.99, samples=80]
    {6*x*(1-x)} \closedcycle;
  \node[font=\small, blue!70!black] at (axis cs:0.5, 2.2)
    {Beta$(2,2)$};
  \node[font=\scriptsize] at (axis cs:0.5, -0.08) {知识不确定};
\end{axis}

% ---- 中图：似然函数 ----
\begin{axis}[
  name=ax2,
  at={(ax1.east)}, anchor=west, xshift=1.8cm,
  width=5.5cm, height=5.5cm,
  xmin=-0.05, xmax=1.05,
  ymin=-0.1, ymax=2.5,
  xlabel={$\theta$},
  ylabel={},
  title={\textbf{似然} $L(\theta\mid x)$},
  axis lines=left,
  xtick={0,0.25,0.5,0.75,1},
  xticklabel style={font=\scriptsize},
  yticklabel style={font=\scriptsize},
  ytick={0,0.5,1.0,1.5,2.0},
  grid=major, grid style={gray!20},
  clip=false
]
  % 似然函数（偏向0.7的钟形，模拟观测数据 7/10 正面）
  % L ∝ theta^7 * (1-theta)^3，峰值约 x=0.7
  \addplot[green!60!black, ultra thick, domain=0.01:0.99, samples=80]
    {2.2*x^7*(1-x)^3/0.0022235};
  \addplot[fill=green!20, opacity=0.5, draw=none, domain=0.01:0.99, samples=80]
    {2.2*x^7*(1-x)^3/0.0022235} \closedcycle;
  \node[font=\small, green!60!black] at (axis cs:0.3, 2.2)
    {$L(\theta\mid x)$};
  \node[font=\scriptsize] at (axis cs:0.5, -0.08) {数据主导};
\end{axis}

% ---- 右图：后验分布 ----
\begin{axis}[
  name=ax3,
  at={(ax2.east)}, anchor=west, xshift=1.8cm,
  width=5.5cm, height=5.5cm,
  xmin=-0.05, xmax=1.05,
  ymin=-0.1, ymax=2.5,
  xlabel={$\theta$},
  ylabel={},
  title={\textbf{后验} $\pi(\theta\mid x)$},
  axis lines=left,
  xtick={0,0.25,0.5,0.75,1},
  xticklabel style={font=\scriptsize},
  yticklabel style={font=\scriptsize},
  ytick={0,0.5,1.0,1.5,2.0},
  grid=major, grid style={gray!20},
  clip=false
]
  % 后验 Beta(9,5)：f ∝ theta^8*(1-theta)^4，峰值在 theta=8/12≈0.667
  \addplot[red!70!black, ultra thick, domain=0.01:0.99, samples=80]
    {2.5*x^8*(1-x)^4/0.000866};
  \addplot[fill=red!20, opacity=0.5, draw=none, domain=0.01:0.99, samples=80]
    {2.5*x^8*(1-x)^4/0.000866} \closedcycle;
  \node[font=\small, red!70!black] at (axis cs:0.3, 2.2)
    {Beta$(9,5)$};
  \node[font=\scriptsize] at (axis cs:0.5, -0.08) {先验与数据融合};
\end{axis}

% ---- 连接箭头 ----
\draw[-{Stealth[length=8pt,width=6pt]}, very thick, orange!80!black]
  ($(ax1.east)!0.5!(ax2.west)$) -- ++(0.5cm,0);
\node[font=\small, orange!80!black, above] at ($(ax1.east)!0.5!(ax2.west)+(0.25cm,0.1cm)$)
  {$\times$};

\draw[-{Stealth[length=8pt,width=6pt]}, very thick, orange!80!black]
  ($(ax2.east)!0.5!(ax3.west)$) -- ++(0.5cm,0);
\node[font=\small, orange!80!black, above] at ($(ax2.east)!0.5!(ax3.west)+(0.25cm,0.1cm)$)
  {$\propto$};

% 贝叶斯公式注释
\node[draw, fill=yellow!15, rounded corners=4pt, font=\small, align=center]
  at ($(ax1.south)!0.5!(ax3.south)+(0,-0.9cm)$)
  {贝叶斯定理：$\pi(\theta\mid x) \propto L(\theta\mid x)\cdot\pi(\theta)$};

% 总标题
\node[font=\large\bfseries] at ($(ax1.north)!0.5!(ax3.north)+(0,1.2cm)$)
  {贝叶斯推断三件套};

\end{tikzpicture}
\end{document}
